\clearpage
\begin{summarybox}
\textbf{\textcolor{CSummary}{一句话核心：特定位置的圆，被抛物线上的三条弦切成“闭合”的三角形。}}

设 $C:y^2=kx$、$\Gamma:(x-h)^2+y^2=r^2$，且 $k,r>0$。
\[
\boxed{h=r+\frac{r^2}{k}},\qquad
\boxed{A_1A_2,A_1A_3\text{ 切圆}\ \Rightarrow\ A_2A_3\text{ 切圆}}.
\]
\textbf{速查两个基础式：}
\[
P(t)=\left(\frac{t^2}{k},t\right),\qquad
\boxed{P(u)P(v):kx-(u+v)y+uv=0},
\]
\[
\boxed{(kh+uv)^2=r^2[k^2+(u+v)^2]}.
\]

\medskip
\textbf{\textcolor{CSummary}{三秒识别流程}}
\[
\boxed{\text{抛物线上三不同点}}\ \longrightarrow\
\boxed{\text{共顶点两弦切圆}}\ \longrightarrow\
\boxed{\text{核对 }h=r+r^2/k}.
\]

\begin{itemize}[leftmargin=*,itemsep=0.42em]
 \item \textbf{直接应用：}满足特殊圆心条件，就知道第三条\emph{直线}也切圆。
 \item \textbf{硬算证明：}两切线方程\textbf{相减求 $b+c$}，再\textbf{相加求 $bc$}；平方和改写成 $(b+c)^2-2bc$。
 \item \textbf{最快证明：}固定 $a$，把 $b,c$ 看作同一个二次方程的两根，用韦达定理；这不是硬算的前置要求。
 \item \textbf{两条边界：}三个点互不相同；说的是所在\emph{直线}切圆，不能自动升级为线段切圆或内切圆。
 \item \textbf{另一位置：} $h=r^2/k-r$ 也可得三直线相切，但要特别检查切点的位置。
\end{itemize}

\medskip
\textbf{\textcolor{CSummary}{真正可迁移的思路：}}当两条共顶点的弦满足同一种条件，题目却问第三条弦时，\textbf{先看第三条弦是否只需要另外两个参数的和与积。} 发现和与积目标后，自然检查“相减 + 相加”，不要一上来就硬解两个点的坐标。
\end{summarybox}
